\chapter*{Phiếu Phản Tư 12 Tuần}
\addcontentsline{toc}{chapter}{Phiếu Phản Tư 12 Tuần}
\markboth{Phiếu Phản Tư 12 Tuần}{Phiếu Phản Tư 12 Tuần}

\begin{truyen}{Dành Cho Học Sinh, Giáo Viên Và Phụ Huynh}{Twelve Weeks of Reflection}
\chuhoa{N}{ếu Quyển XXXIII chỉ dừng} ở việc đọc xong rồi gấp lại, tác dụng của nó sẽ ngắn. Phần này biến nội dung của sách thành 12 tuần phản tư. Mỗi phiếu có một trọng tâm, một câu soi mình, một việc nhỏ phải làm, và một khoảng trống để viết thật ra điều đã đổi hoặc chưa đổi.
\textit{(If Volume XXXIII ends at being read and closed, its effect will be short. This section turns the book into twelve weeks of reflection. Each sheet has a focus, a self-question, one small action, and space to record what changed or failed to change.)}
\end{truyen}

\newcommand{\phieuphan}[4]{%
  \begin{tcolorbox}[
    colback=white,
    colframe=bookprimary!55,
    arc=4pt,
    boxrule=0.8pt,
    title={#1},
    fonttitle=\bfseries,
    breakable,
    enhanced
  ]
  \textbf{Trọng tâm:} #2

  \vspace{0.2cm}
  \textbf{Tự hỏi:} #3

  \vspace{0.2cm}
  \textbf{Việc nhỏ phải làm trong tuần:} #4

  \vspace{0.3cm}
  \textit{Điều tôi đã làm được:}

  \vspace{0.25cm}
  \hrulefill

  \vspace{0.35cm}
  \hrulefill

  \vspace{0.35cm}
  \textit{Điều tôi còn chống chế hoặc chưa vượt qua:}

  \vspace{0.25cm}
  \hrulefill

  \vspace{0.35cm}
  \hrulefill

  \vspace{0.35cm}
  \textit{Một người có thể nhắc tôi giữ lời tuần này: \hrulefill}
  \end{tcolorbox}
  \vspace{0.45cm}
}

\phieuphan{Tuần 1. Em Đang Mệt Hay Đang Né?}{Gọi đúng tên áp lực}{Tuần này, điều gì làm tôi mệt thật và điều gì tôi đang phóng đại để xin buông?}{Chọn một việc khó nhất và chia nó thành ba bước nhỏ để bắt đầu.}

\phieuphan{Tuần 2. Bình Yên Hay Dễ Dãi?}{Tự chăm sóc và tự chiều mình}{Tôi đang nghỉ để hồi phục, hay đang nghỉ để đỡ phải lớn lên?}{Bỏ một thói quen ru ngủ mình và thay bằng một nhịp sinh hoạt rõ ràng hơn.}

\phieuphan{Tuần 3. Một Chữ Không}{Áp lực đồng trang}{Tôi đã đồng ý điều gì chỉ vì sợ lạc nhóm?}{Tập nói ``không'' với một lời rủ rê không hợp lý, dù rất nhỏ.}

\phieuphan{Tuần 4. Nhìn Người Lớn Công Bằng}{Công bằng với lỗi của mọi phía}{Khi người lớn sai, tôi có đang muốn công bằng hay chỉ muốn hả hê?}{Viết lại một tình huống mà cả tôi và người lớn đều có phần phải sửa.}

\phieuphan{Tuần 5. Phản Biện Có Trí Tuệ}{Nội quy và cách đặt câu hỏi}{Tôi đang phản biện để làm rõ điều đúng, hay phản ứng để thắng miệng?}{Chọn một điều tôi thấy vô lý và viết lại thành ba câu hỏi tử tế.}

\phieuphan{Tuần 6. Xin Lỗi Không Mặc Cả}{Nhân cách và chuẩn mực}{Lần gần đây nhất tôi xin lỗi, tôi xin lỗi để sửa hay để thoát?}{Xin lỗi một việc cụ thể mà không kèm chống chế phía sau.}

\phieuphan{Tuần 7. Bênh Nhưng Không Che}{Tình thương có tỉnh táo}{Khi thương một người, tôi có đang giúp họ trưởng thành hay giúp họ né hậu quả?}{Nói một câu góp ý thật lòng với người mình thương theo cách không làm nhục họ.}

\phieuphan{Tuần 8. Thước Đo Của Riêng Mình}{So sánh và tự giá}{Câu so sánh nào vẫn còn làm tôi thấy mình thua kém?}{Tự đặt một thước đo tiến bộ cá nhân trong đúng một lĩnh vực.}

\phieuphan{Tuần 9. Từ Giận Sang Kế Hoạch}{Đối thoại ba bên}{Khi có vấn đề, tôi giỏi trách hơn hay giỏi nghĩ bước tiếp theo?}{Biến một lời than của mình thành một kế hoạch nhỏ có hạn chót.}

\phieuphan{Tuần 10. Ước Mơ Cần Nền}{Ảo tưởng thành công}{Tôi đang mê con đường nào vì hiểu mình, và mê con đường nào vì nghe nó ngầu?}{Ghi ra ba phần nền tôi phải rèn nếu vẫn muốn theo ước mơ ấy.}

\phieuphan{Tuần 11. Trải Nghiệm Phải Để Lại Gì}{Cảm giác mới lạ và sự lớn lên}{Những trải nghiệm gần đây đã làm tôi sâu hơn ở chỗ nào?}{Sau một hoạt động hoặc chuyến đi, viết lại đúng một kỹ năng hoặc một thay đổi nhận thức còn giữ được.}

\phieuphan{Tuần 12. Thái Độ Đi Vào Đời}{Lao động, kỷ luật, phản hồi}{Nếu ngày mai đi làm, lý sự nào của tôi sẽ vỡ đầu tiên?}{Tập đúng giờ hơn, phản hồi trưởng thành hơn, và chịu được một góp ý khó nghe mà không quy thành bị ghét.}

\begin{baihoc}
Mười hai phiếu này không cần làm thật đẹp. Chúng chỉ cần được điền thật. Nhiều khi một câu trả lời thành thật nhưng xấu xí còn có ích hơn mười câu rất hay nhưng viết để trình diễn.
\textit{(These twelve sheets do not need to be completed beautifully. They only need to be completed honestly. One ugly but truthful answer is often more useful than ten elegant answers written for display.)}
\end{baihoc}

\ngancach